\chapter{大规模算子分裂算法}

\label{chap:10}

第 \ref{chap:09} 章已经把凸最小化、变分不等式与 KKT 系统统一改写成单调包含
\[
0\in Au.
\]
因此，本章不再重新发明新的最优性条件，而是讨论怎样围绕 resolvent、prox 与固定点结构设计可扩展的迭代法。所谓“大规模”，在这里并不只是维数高，而是指目标函数、约束与对偶变量往往天然分裂成若干模块，直接求解完整 KKT 系统代价过高；真正可行的方案，是只调用梯度、投影、prox 或局部线性算子，并把耦合留给迭代去吸收。

本章沿四条主线展开。第 \ref{sec:10-1} 节讨论前向--后向分裂，把“一个可显式前向步的单值算子”和“一个可通过 resolvent 处理的极大单调算子”组合起来；这是第 \ref{sec:9-2} 节 prox 和第 \ref{sec:9-4} 节 zero finding 模板的最直接算法化。第 \ref{sec:10-2} 节引入 Douglas--Rachford 分裂，说明反射算子如何把两个极大单调算子的和改写成固定点问题。第 \ref{sec:10-3} 节随后把增广拉格朗日与 ADMM 挂回第 \ref{sec:7-4} 节的 KKT 语言，解释 Boyd 书中最常见的原--对偶迭代其实就是 Douglas--Rachford 在对偶侧的变量消去形式。最后，第 \ref{sec:10-4} 节讨论惰性与加速模板，说明 Hilbert 空间中的加速为什么必须借助能量函数与稳定性分析，而不能把有限维口号不加区分地直接搬到无穷维情形。

本章的正文主舞台是 Hilbert 空间。原因并不神秘：第 \ref{sec:3-1} 节定理~\ref{thm:riesz-representation-hilbert} 使得对偶元素能够被向量化表示，而第 \ref{sec:9-2} 节的 prox、第 \ref{sec:9-3} 节的 Moreau--Yosida 平滑以及平均映射的非扩张几何都在这一背景下最为透明。更一般的 Banach 空间版本当然存在，但其中投影、cocoercivity、firm nonexpansive 与反射算子的几何优势会明显减弱，因此本章只在 remark 中点到这些边界，而不把正文写成一部一般 Banach 分裂算法专著。

从叙事上看，本章恰好是第 \ref{chap:11} 章数值求解与第 \ref{chap:12} 章分布式实现之间的桥。前者关心连续时间极限、离散稳定性与收敛轨道，后者关心共识、通信与模块化实现；而它们共同依赖的，是本章建立的“算子分裂 = 固定点求解 = 局部可计算更新”这一接口。Boyd 的有限维算法表面上看像一组彼此分离的工程配方，本章的任务则是说明，它们其实共享同一条单调算子主线。

\section{前向--后向分裂算法}

\label{sec:10-1}

第 \ref{sec:9-4} 节把极小化与 KKT 系统改写成单调包含之后，最先出现的问题是：若算子 $A$ 还可以拆成“容易前向评估的一部分”和“容易做 resolvent 的一部分”，能否不必在每一步都求解完整隐式方程？前向--后向分裂正是对这个问题的回答。它在实践里常被直接叫作前向-后向算法；当“前向部分”来自光滑凸函数梯度，而“后向部分”来自非光滑正则项次微分时，它就退化为最熟悉的 prox-gradient 迭代。

\subsection{问题模板与基本定义}

\begin{definition}[cocoercive 算子]\label{def:cocoercive-operator}
设 $H$ 为 Hilbert 空间，$B\colon H\to H$ 为单值算子，$\beta>0$。若对任意 $x,y\in H$ 都有
\[
\langle Bx-By,x-y\rangle \ge \beta \|Bx-By\|^2,
\]
则称 $B$ 是 $\beta$-\emph{cocoercive} 的。
\end{definition}

\begin{definition}[前向--后向迭代]\label{def:forward-backward-iteration}
设 $H$ 为 Hilbert 空间，$B\colon H\to H$ 为单值算子，$C\colon H\rightrightarrows H$ 为极大单调算子，$\gamma>0$。定义
\[
T_\gamma:=J_{\gamma C}(I-\gamma B).
\]
给定初值 $x^0\in H$，迭代
\[
x^{k+1}=T_\gamma x^k=J_{\gamma C}(x^k-\gamma Bx^k),\qquad k=0,1,2,\dots
\]
称为求解
\[
0\in Bx+Cx
\]
的\emph{前向--后向迭代}。
\end{definition}

\begin{remark}
定义~\ref{def:forward-backward-iteration} 之所以自然，是因为它严格沿着第 \ref{sec:9-4} 节推论~\ref{cor:proximal-point-template} 的接口来设计：$C$ 仍然通过 resolvent 处理，但 $B$ 不再被整体塞进隐式步，而是只做一次显式前向评估。若 $C=\partial g$，则由命题~\ref{prop:prox-resolvent-subdifferential}，后向步就是 prox；若 $B=\nabla h$，则前向步就是普通梯度下降。因此它同时连接了第 \ref{chap:05} 章的微分语言与第 \ref{chap:09} 章的单调算子语言。
\end{remark}

\subsection{固定点与零点的对应}

\begin{proposition}[前向--后向固定点表述]\label{prop:forward-backward-fixed-point}
设 $B$ 与 $C$ 如定义~\ref{def:forward-backward-iteration}，$\gamma>0$。则
\[
\operatorname{Fix}T_\gamma=\operatorname{zer}(B+C).
\]
换言之，对任意 $x\in H$，
\[
x=T_\gamma x
\quad\Longleftrightarrow\quad
0\in Bx+Cx.
\]
\end{proposition}

\begin{proof}
由 resolvent 的定义，
\[
x=T_\gamma x
\quad\Longleftrightarrow\quad
x=J_{\gamma C}(x-\gamma Bx)
\quad\Longleftrightarrow\quad
x-\gamma Bx\in x+\gamma Cx.
\]
最后一个关系等价于
\[
-Bx\in Cx,
\]
亦即
\[
0\in Bx+Cx.
\]
反向推理同理成立。
\end{proof}

\subsection{平均映射结构与收敛}

\begin{theorem}[前向--后向分裂的收敛定理]\label{thm:forward-backward-convergence}
设 $H$ 为 Hilbert 空间，$C\colon H\rightrightarrows H$ 为极大单调算子，$B\colon H\to H$ 为 $\beta$-cocoercive 单值算子，并假设
\[
\operatorname{zer}(B+C)\neq \varnothing.
\]
若步长满足
\[
0<\gamma<2\beta,
\]
则由定义~\ref{def:forward-backward-iteration} 生成的序列 $\{x^k\}$ 弱收敛到某个解 $\bar x\in \operatorname{zer}(B+C)$。
\end{theorem}

\begin{proof}
先说明 $T_\gamma$ 是平均映射。令
\[
\alpha:=\frac{\gamma}{2\beta}\in(0,1).
\]
由定义~\ref{def:cocoercive-operator}，
\[
\|(I-2\beta B)x-(I-2\beta B)y\|^2
=
\|x-y\|^2-4\beta\langle Bx-By,x-y\rangle+4\beta^2\|Bx-By\|^2
\le \|x-y\|^2.
\]
故 $N:=I-2\beta B$ 非扩张，而
\[
I-\gamma B=(1-\alpha)I+\alpha N
\]
是 $\alpha$-averaged 映射。另一方面，由第 \ref{sec:9-1} 节定理~\ref{thm:minty-surjectivity}，$J_{\gamma C}$ 处处有定义；再由第 \ref{sec:9-2} 节命题~\ref{prop:prox-firm-nonexpansive} 在 resolvent 语言下的对应形式，$J_{\gamma C}$ 是 firmly nonexpansive，也就是 $1/2$-averaged。平均映射的复合仍是平均映射，因此存在某个 $\theta\in(0,1)$ 使
\[
T_\gamma=(1-\theta)I+\theta N_\gamma
\]
其中 $N_\gamma$ 非扩张。

再由命题~\ref{prop:forward-backward-fixed-point}，
\[
\operatorname{Fix}T_\gamma=\operatorname{zer}(B+C)\neq\varnothing.
\]
任取 $z\in \operatorname{Fix}T_\gamma$。平均映射的标准不等式给出
\[
\|T_\gamma x-z\|^2
\le
\|x-z\|^2-\frac{1-\theta}{\theta}\|(I-T_\gamma)x\|^2.
\]
将 $x=x^k$ 代入并使用 $x^{k+1}=T_\gamma x^k$，得到
\[
\|x^{k+1}-z\|^2
\le
\|x^k-z\|^2-\frac{1-\theta}{\theta}\|x^k-x^{k+1}\|^2.
\]
因此 $\{x^k\}$ 对解集是 Fej\'er 单调的，特别地有界，并且
\[
\sum_{k=0}^\infty \|x^{k+1}-x^k\|^2<\infty.
\]
于是 $\|x^{k+1}-x^k\|\to 0$，即序列渐近正则。

设 $\bar x$ 为 $\{x^k\}$ 的任意弱聚点，则存在子列 $x^{k_j}\rightharpoonup \bar x$。由于 $x^{k_j}-T_\gamma x^{k_j}\to 0$，而平均映射满足 demiclosedness 原理，故
\[
\bar x\in \operatorname{Fix}T_\gamma.
\]
再由 Hilbert 空间中的 Opial 引理，整个序列弱收敛到某个 $\bar x\in \operatorname{Fix}T_\gamma=\operatorname{zer}(B+C)$。
\end{proof}

\begin{corollary}[prox-gradient 是前向--后向分裂的特例]\label{cor:prox-gradient-specialization}
设 $h\colon H\to \mathbb R$ 为凸且 Fr\'echet 可微函数，且 $\nabla h$ 为 $\beta$-cocoercive；设 $g\colon H\to \mathbb R\cup\{+\infty\}$ 为 proper、凸且下半连续的泛函。则对问题
\[
\min_{x\in H}\{h(x)+g(x)\},
\]
前向--后向迭代退化为
\[
x^{k+1}
=
\operatorname{prox}_{\gamma g}(x^k-\gamma \nabla h(x^k)),
\qquad 0<\gamma<2\beta.
\]
这正是 Hilbert 空间中的 prox-gradient 算法。
\end{corollary}

\begin{proof}
把最优性条件写成
\[
0\in \nabla h(x)+\partial g(x).
\]
在定义~\ref{def:forward-backward-iteration} 中取 $B=\nabla h$、$C=\partial g$。由命题~\ref{prop:prox-resolvent-subdifferential}，
\[
J_{\gamma\partial g}=\operatorname{prox}_{\gamma g}.
\]
因此前向--后向迭代恰好写成所述的 prox-gradient 形式。收敛性则由定理~\ref{thm:forward-backward-convergence} 直接得到。
\end{proof}

\begin{remark}
在具体建模时，常把 $h$ 写成数据拟合项、把 $g$ 写成正则项。若拟合项先经过线性算子复合，则第 \ref{sec:5-4} 节命题~\ref{prop:subdifferential-linear-pullback} 说明其最优性系统依旧能嵌入定义~\ref{def:forward-backward-iteration}。这就是为什么前向--后向分裂能自然覆盖经验风险最小化、逆问题与变分正则化模型。
\end{remark}

\subsection{典型例子}

\begin{example}[$f+g$ 型凸最小化]
对
\[
\min_x \{h(x)+g(x)\},
\]
若 $h$ 平滑且 $\nabla h$ cocoercive，$g$ 允许非光滑，则推论~\ref{cor:prox-gradient-specialization} 给出标准 prox-gradient 更新。把这个例子放在这里，是为了强调：前向--后向分裂不是另起炉灶的新算法，而是第 \ref{sec:9-2} 节 prox 与普通梯度步的自然拼接。
\end{example}

\begin{example}[Lasso 的 ISTA]
在 $\mathbb R^n$ 上考虑
\[
\min_x \frac{1}{2}\|Ax-b\|^2+\lambda \|x\|_1.
\]
此时 $h(x)=\frac12\|Ax-b\|^2$，$g(x)=\lambda\|x\|_1$。若 $\|A\|^2\le L$，则 $\nabla h$ 是 $1/L$-cocoercive，于是前向-后向迭代成为
\[
x^{k+1}
=
\operatorname{prox}_{\gamma\lambda\|\cdot\|_1}
\bigl(x^k-\gamma A^\ast(Ax^k-b)\bigr),
\qquad 0<\gamma<\frac{2}{L}.
\]
这就是 Boyd 书中最常见的 ISTA 形式。
\end{example}

\begin{example}[带正则项的经验风险最小化]
在 Hilbert 空间中，若
\[
h(x)=\frac{1}{m}\sum_{i=1}^m \ell_i(\langle a_i,x\rangle),
\qquad
g(x)=\tau \|x\|_1 \ \text{或}\ \tau\|x\|_\ast,
\]
则前向步只需评估梯度，后向步只需调用稀疏或低秩 prox。把这个例子放在这里，是为了说明“大规模”并不要求每一步求精确解，而是要求每一步只做局部可扩展的运算。
\end{example}

\subsection*{有限维对照}

Boyd 往往把 prox-gradient 画成“先梯度下降，再做阈值化”的两步法。本节说明，这种 Euclidean 图像只是推论~\ref{cor:prox-gradient-specialization} 的有限维外观。真正起作用的结构是：$\nabla h$ 的 cocoercive 性保证前向步可控，resolvent 的 firm nonexpansive 性保证后向步稳定，而定理~\ref{thm:forward-backward-convergence} 则把两者合成为一个平均映射的收敛论证。有限维不过是这一 Hilbert 结构的最直观坍缩。

\subsection*{习题（待补）}

\begin{exercise}[平均映射验证占位]
补一题：在定理~\ref{thm:forward-backward-convergence} 的假设下，直接验证 $I-\gamma B$ 为 averaged 映射，并据此推出 $T_\gamma$ 的 averaged 性。
\end{exercise}

\begin{exercise}[ISTA 桥接题占位]
补一题：对 Lasso 模型写出前向-后向迭代，并利用例中的软阈值公式把它化为显式坐标更新。
\end{exercise}


\section{Douglas--Rachford 分裂}

\label{sec:10-2}

前向-后向分裂适合处理“一个显式、一个隐式”的和；若两部分都是极大单调算子，直接做前向步往往并不现实。这时更合适的对象不是梯度步，而是反射算子与其平均。Douglas-Rachford 方法正是在这一点上取代了显式梯度：它不去近似 $A+B$，而是把两个 resolvent 组织成一个固定点映射，然后再从固定点中恢复零点。

\subsection{反射算子与迭代定义}

\begin{definition}[反射算子]\label{def:reflection-operator}
设 $H$ 为 Hilbert 空间，$A\colon H\rightrightarrows H$ 为极大单调算子，$\gamma>0$。定义
\[
R_{\gamma A}:=2J_{\gamma A}-I.
\]
该映射称为算子 $A$ 的 \emph{反射算子}。
\end{definition}

\begin{definition}[Douglas--Rachford 迭代]\label{def:douglas-rachford-iteration}
设 $A,B\colon H\rightrightarrows H$ 为极大单调算子，$\gamma>0$。定义
\[
T_{\mathrm{DR}}
:=
\frac12\bigl(I+R_{\gamma A}R_{\gamma B}\bigr).
\]
给定初值 $z^0\in H$，迭代
\[
z^{k+1}=T_{\mathrm{DR}}z^k
\]
称为求解
\[
0\in Ax+Bx
\]
的 \emph{Douglas--Rachford 分裂}。
\end{definition}

\begin{remark}
由第 \ref{sec:9-1} 节定理~\ref{thm:minty-surjectivity}，$J_{\gamma A}$ 与 $J_{\gamma B}$ 处处有定义，因此定义~\ref{def:reflection-operator} 和定义~\ref{def:douglas-rachford-iteration} 都是良定的。与第 \ref{sec:10-1} 节不同，这里不要求某一部分是显式可微项；两个块都只通过 resolvent 进入算法。
\end{remark}

\subsection{固定点与零点恢复}

\begin{proposition}[Douglas--Rachford 固定点与零点]\label{prop:dr-fixed-point-zer-sum}
设 $A,B$ 为极大单调算子，$\gamma>0$。则有：
\begin{enumerate}
  \item 若 $z\in \operatorname{Fix}T_{\mathrm{DR}}$，并令
  \[
  x:=J_{\gamma B}z,
  \]
  则 $x\in \operatorname{zer}(A+B)$；
  \item 反过来，若 $x\in \operatorname{zer}(A+B)$，则存在 $z\in \operatorname{Fix}T_{\mathrm{DR}}$ 使 $x=J_{\gamma B}z$。
\end{enumerate}
\end{proposition}

\begin{proof}
先证第一条。设 $z=T_{\mathrm{DR}}z$，即
\[
z=\frac12(z+R_{\gamma A}R_{\gamma B}z),
\]
从而
\[
z=R_{\gamma A}R_{\gamma B}z.
\]
记
\[
x:=J_{\gamma B}z,
\qquad
R_{\gamma B}z=2x-z.
\]
由上式和反射的定义，
\[
z
=
2J_{\gamma A}(R_{\gamma B}z)-R_{\gamma B}z,
\]
因此
\[
J_{\gamma A}(R_{\gamma B}z)
=
\frac12(z+R_{\gamma B}z)
=
x.
\]
由 resolvent 的定义，
\[
z-x\in \gamma Bx,
\qquad
R_{\gamma B}z-x=x-z\in \gamma A x.
\]
两式相加得到
\[
0\in \gamma A x+\gamma Bx,
\]
即
\[
0\in Ax+Bx.
\]

再证第二条。若 $x\in \operatorname{zer}(A+B)$，则可取
\[
a\in Ax,\qquad b\in Bx,\qquad a+b=0.
\]
令
\[
z:=x+\gamma b.
\]
则
\[
z-x\in \gamma Bx,
\]
故 $x=J_{\gamma B}z$。同时
\[
R_{\gamma B}z=2x-z=x-\gamma b=x+\gamma a,
\]
于是
\[
R_{\gamma B}z-x\in \gamma A x,
\]
故 $x=J_{\gamma A}(R_{\gamma B}z)$。代回反射定义可得
\[
R_{\gamma A}R_{\gamma B}z
=
2J_{\gamma A}(R_{\gamma B}z)-R_{\gamma B}z
=
2x-(x-\gamma b)
=
x+\gamma b
=
z.
\]
于是 $z\in \operatorname{Fix}T_{\mathrm{DR}}$。
\end{proof}

\subsection{平均映射性质}

\begin{theorem}[Douglas--Rachford 映射是 averaged]\label{thm:douglas-rachford-averaged}
设 $A,B\colon H\rightrightarrows H$ 为极大单调算子，$\gamma>0$，并假设
\[
\operatorname{zer}(A+B)\neq \varnothing.
\]
则 $T_{\mathrm{DR}}$ 是 firmly nonexpansive，从而尤其是 averaged 映射。由定义~\ref{def:douglas-rachford-iteration} 生成的序列 $\{z^k\}$ 弱收敛到某个 $\bar z\in \operatorname{Fix}T_{\mathrm{DR}}$，且影子序列
\[
x^k:=J_{\gamma B}z^k
\]
弱收敛到某个 $\bar x\in \operatorname{zer}(A+B)$。
\end{theorem}

\begin{proof}
由 resolvent 的 firm nonexpansive 性，$J_{\gamma A}$ 与 $J_{\gamma B}$ 都是 firmly nonexpansive。于是对任意 $u,v\in H$，
\[
\|R_{\gamma A}u-R_{\gamma A}v\|
=
\|2J_{\gamma A}u-u-(2J_{\gamma A}v-v)\|
\le \|u-v\|.
\]
故 $R_{\gamma A}$ 非扩张；同理 $R_{\gamma B}$ 也非扩张，于是复合 $R_{\gamma A}R_{\gamma B}$ 仍非扩张。由定义~\ref{def:douglas-rachford-iteration}，
\[
T_{\mathrm{DR}}=\frac12 I+\frac12 R_{\gamma A}R_{\gamma B}
\]
是 $1/2$-averaged，即 firmly nonexpansive。

再由命题~\ref{prop:dr-fixed-point-zer-sum}，$\operatorname{zer}(A+B)\neq\varnothing$ 蕴含 $\operatorname{Fix}T_{\mathrm{DR}}\neq\varnothing$。于是可像定理~\ref{thm:forward-backward-convergence} 那样，对 averaged 映射使用 Fej\'er 单调与 Opial 引理，得到 $z^k\rightharpoonup \bar z\in \operatorname{Fix}T_{\mathrm{DR}}$。

由于 resolvent 非扩张，$\{x^k\}$ 有界。若 $x^{k_j}\rightharpoonup \bar x$，则由
\[
x^{k_j}=J_{\gamma B}z^{k_j},
\qquad
z^{k_j}\rightharpoonup \bar z
\]
以及 resolvent 图像的闭性可知 $\bar x=J_{\gamma B}\bar z$。再由命题~\ref{prop:dr-fixed-point-zer-sum}，$\bar x\in \operatorname{zer}(A+B)$。因此整个影子序列也弱收敛到某个解。
\end{proof}

\begin{corollary}[可行性交是 Douglas--Rachford 的特例]\label{cor:dr-feasibility-special-case}
设 $C,D\subset H$ 为非空闭凸集，并考虑可行性交问题
\[
\text{求 }x\in C\cap D.
\]
令
\[
A=N_C,\qquad B=N_D.
\]
则定义~\ref{def:douglas-rachford-iteration} 退化为
\[
T_{\mathrm{DR}}
=
\frac12\bigl(I+(2P_C-I)(2P_D-I)\bigr),
\]
而任意固定点 $\bar z$ 的影子点
\[
\bar x=P_D\bar z
\]
属于 $C\cap D$。
\end{corollary}

\begin{proof}
由第 \ref{sec:5-4} 节推论~\ref{cor:normal-cone-indicator} 与第 \ref{sec:9-2} 节投影例子，
\[
J_{\gamma N_C}=P_C,\qquad J_{\gamma N_D}=P_D.
\]
于是反射算子写成
\[
R_{\gamma N_C}=2P_C-I,\qquad R_{\gamma N_D}=2P_D-I.
\]
其余结论由定理~\ref{thm:douglas-rachford-averaged} 和命题~\ref{prop:dr-fixed-point-zer-sum} 直接得到。
\end{proof}

\begin{remark}
第 \ref{sec:10-2} 节的意义不只在于可行性交。第 \ref{sec:10-3} 节会看到，ADMM 并不是一套与 Douglas-Rachford 平行的经验算法；恰恰相反，ADMM 可以被看成 Douglas--Rachford 在对偶侧的变量消去形式。也正因为如此，理解反射算子与固定点结构，是理解 ADMM 的最快途径。
\end{remark}

\subsection{典型例子}

\begin{example}[两个闭凸集的可行性交]
在例子最基本的情形中，$C$ 与 $D$ 分别代表两个独立约束模块，例如一个数据一致性集合和一个结构先验集合。Douglas-Rachford 的每一步都只需做两次投影与两次反射，却能把可行性交改写成单个固定点问题。这种“把复杂交集拆成简单块”的思想，是大规模模型中最关键的算法动机之一。
\end{example}

\begin{example}[欧氏空间中的双反射]
在 $\mathbb R^n$ 中，当 $C$ 与 $D$ 是仿射子空间时，$2P_C-I$ 与 $2P_D-I$ 就是熟悉的镜面对称。把这个例子放在这里，是为了说明 Douglas-Rachford 的几何直觉确实可以在有限维里画出来；但真正支撑它的不是图形，而是定理~\ref{thm:douglas-rachford-averaged} 的平均映射结构。
\end{example}

\begin{example}[图像恢复中的一致性约束]
在图像恢复或逆问题中，常把“符合观测”的约束与“满足先验”的约束分开建模，再对两者的交施加 Douglas-Rachford 分裂。把这个例子放在这里，是为了说明可行性交并不只是几何练习，而是后续 ADMM、投影算法和一致性恢复模型的共同母语。
\end{example}

\subsection*{有限维对照}

Boyd 书里常把 Douglas-Rachford 与交替投影、反射法并列介绍，读起来像几种彼此不同的几何技巧。本节的结论是：它们共享同一个固定点框架。有限维里，反射算子可以直接画成镜像，因而算法看起来很直观；无穷维里，这种图像直觉不再可靠，但定理~\ref{thm:douglas-rachford-averaged} 和命题~\ref{prop:dr-fixed-point-zer-sum} 依然成立，因此 Boyd 中的几何配方其实只是更一般单调算子结构的坍缩版本。

\subsection*{习题（待补）}

\begin{exercise}[反射算子占位]
补一题：证明对非空闭凸集 $C$，反射算子 $2P_C-I$ 非扩张，并据此解释定义~\ref{def:reflection-operator} 在法锥情形中的几何意义。
\end{exercise}

\begin{exercise}[可行性交桥接题占位]
补一题：把一个具体的可行性交模型改写为 $0\in N_Cx+N_Dx$，并写出对应的 Douglas-Rachford 迭代。
\end{exercise}


\section{交替方向乘子法的泛函视角}

\label{sec:10-3}

第 \ref{sec:7-4} 节已经说明，带约束凸优化的真正核心对象不是原问题本身，而是一个 KKT 原--对偶系统。ADMM 的价值就在于，它并不试图一次性解这个系统，而是借助增广拉格朗日把耦合项放进一个二次罚项里，再通过交替最小化与乘子更新逐步逼近鞍点。有限维里，这常被当成 Boyd 式工程算法来讲；本节要说明的是，它其实可以完全挂回第 \ref{sec:10-2} 节的 Douglas--Rachford 分裂与第 \ref{sec:7-2} 节的标准型接口。

\subsection{两块模型与增广拉格朗日}

\begin{definition}[ADMM 的增广拉格朗日]\label{def:augmented-lagrangian-admm}
设 $X,Z,Y$ 为 Hilbert 空间，$A\colon X\to Y$、$B\colon Z\to Y$ 为有界线性算子，$b\in Y$。考虑两块凸问题
\[
\min_{x\in X,\ z\in Z}\{f(x)+g(z):Ax+Bz=b\},
\]
其中 $f\colon X\to \mathbb R\cup\{+\infty\}$、$g\colon Z\to \mathbb R\cup\{+\infty\}$ 为 proper、凸且下半连续泛函。对给定罚参数 $\rho>0$，定义其\emph{增广拉格朗日}为
\[
\mathcal L_\rho(x,z;y)
:=
f(x)+g(z)+\langle y,Ax+Bz-b\rangle+\frac{\rho}{2}\|Ax+Bz-b\|^2.
\]
\end{definition}

\begin{definition}[ADMM 迭代]\label{def:admm-iteration}
在定义~\ref{def:augmented-lagrangian-admm} 的条件下，给定 $(x^0,z^0,y^0)$ 与 $\rho>0$，定义
\[
x^{k+1}\in \arg\min_x \mathcal L_\rho(x,z^k;y^k),
\]
\[
z^{k+1}\in \arg\min_z \mathcal L_\rho(x^{k+1},z;y^k),
\]
\[
y^{k+1}=y^k+\rho(Ax^{k+1}+Bz^{k+1}-b).
\]
这组更新称为 \emph{ADMM}。相应地，定义原残差与对偶残差
\[
r^{k+1}:=Ax^{k+1}+Bz^{k+1}-b,
\qquad
s^{k+1}:=\rho B^\ast(z^{k+1}-z^k).
\]
\end{definition}

\begin{remark}
定义~\ref{def:admm-iteration} 中的乘子更新，正是把约束违反量 $r^{k+1}$ 回灌到对偶变量中。它与第 \ref{sec:7-4} 节定义~\ref{def:kkt-cone-system} 的关系是直接的：若 $r^{k+1}\to 0$ 且两个块的最优性条件稳定下来，就会逼近 KKT 系统。也因此，ADMM 最终要验证的不是某个经验停止准则，而是原--对偶残差是否把 KKT 条件逼近到可接受精度。
\end{remark}

\subsection{ADMM 与 Douglas--Rachford 的接口}

\begin{proposition}[ADMM 是 Douglas--Rachford 分裂的变量消去形式]\label{prop:admm-as-dr-splitting}
在标准的资格条件下，把定义~\ref{def:augmented-lagrangian-admm} 的原问题经由第 \ref{sec:7-2} 节命题~\ref{prop:abstract-standard-form-to-fenchel} 写成对偶问题后，可得到两个极大单调算子
\[
M_f:=\partial\bigl(f^\ast\circ(-A^\ast)\bigr)+b,
\qquad
M_g:=\partial\bigl(g^\ast\circ(-B^\ast)\bigr).
\]
对偶变量的一个仿射重参数化满足第 \ref{sec:10-2} 节定义~\ref{def:douglas-rachford-iteration} 的 Douglas--Rachford 迭代。因此，ADMM 可以看成对偶侧 DR 分裂在原变量与乘子变量上的显式展开。
\end{proposition}

\begin{proof}
由 $x$-步的最优性条件，
\[
0\in \partial f(x^{k+1})+A^\ast y^k+\rho A^\ast(Ax^{k+1}+Bz^k-b).
\]
再利用乘子更新 $y^{k+1}=y^k+\rho r^{k+1}$，可重写为
\[
-A^\ast\bigl(y^{k+1}+\rho B(z^k-z^{k+1})\bigr)\in \partial f(x^{k+1}).
\]
同样地，$z$-步给出
\[
-B^\ast y^{k+1}\in \partial g(z^{k+1}).
\]
把这两条最优性条件通过 Fenchel 反演改写到对偶空间，并引入标准的中间变量，就会得到两个 resolvent 的交替作用；其组合正是
\[
\frac12\bigl(I+R_{\rho M_f}R_{\rho M_g}\bigr)
\]
的 Douglas--Rachford 形式。这里真正重要的不是坐标层面的繁琐代换，而是结构：ADMM 的“交替最小化 + 乘子更新”并不是额外发明的一套规则，而是第 \ref{sec:10-2} 节 DR 分裂在对偶侧的等价写法。
\end{proof}

\subsection{收敛与 KKT 极限}

\begin{theorem}[ADMM 的原--对偶收敛]\label{thm:admm-primal-dual-convergence}
设定义~\ref{def:augmented-lagrangian-admm} 的问题存在鞍点 $(\bar x,\bar z,\bar y)$，并假设定义~\ref{def:admm-iteration} 生成的序列有界。则原残差与对偶残差满足
\[
r^k\to 0,
\qquad
s^k\to 0.
\]
并且每个弱聚点 $(x^\infty,z^\infty,y^\infty)$ 都满足相应的 KKT 系统；若 KKT 解集在所考虑模型下是单点，则整个序列弱收敛到该原--对偶解。
\end{theorem}

\begin{proof}
先写出 ADMM 两个子问题的一阶最优性条件。对 $x$-步，
\[
0\in \partial f(x^{k+1})+A^\ast y^k+\rho A^\ast(Ax^{k+1}+Bz^k-b),
\]
即
\[
-A^\ast\bigl(y^{k+1}+\rho B(z^k-z^{k+1})\bigr)\in \partial f(x^{k+1}).
\]
对 $z$-步，
\[
0\in \partial g(z^{k+1})+B^\ast y^k+\rho B^\ast(Ax^{k+1}+Bz^{k+1}-b),
\]
于是
\[
-B^\ast y^{k+1}\in \partial g(z^{k+1}).
\]
再与鞍点的 KKT 条件
\[
-A^\ast\bar y\in \partial f(\bar x),
\qquad
-B^\ast\bar y\in \partial g(\bar z),
\qquad
A\bar x+B\bar z-b=0
\]
配对，并利用次微分的单调性，得到
\[
\bigl\langle A(x^{k+1}-\bar x),\, y^{k+1}-\bar y+\rho B(z^k-z^{k+1}) \bigr\rangle \le 0,
\]
\[
\bigl\langle B(z^{k+1}-\bar z),\, y^{k+1}-\bar y \bigr\rangle \le 0.
\]
把两式相加，并用
\[
A(x^{k+1}-\bar x)=r^{k+1}-B(z^{k+1}-\bar z)
\]
代入，可得
\[
\langle r^{k+1},y^{k+1}-\bar y\rangle
\le
\rho \langle B(z^{k+1}-\bar z),\,B(z^k-z^{k+1})\rangle.
\]
再利用恒等式
\[
y^{k+1}-y^k=\rho r^{k+1}
\]
以及极化公式
\[
2\langle u-v,u-w\rangle
=
\|u-w\|^2-\|v-w\|^2+\|u-v\|^2,
\]
可将上式化为
\[
\|y^{k+1}-\bar y\|^2+\rho^2\|B(z^{k+1}-\bar z)\|^2
\le
\|y^k-\bar y\|^2+\rho^2\|B(z^k-\bar z)\|^2
\]
\[
\qquad\qquad
-\rho^2\|r^{k+1}\|^2-\rho^2\|B(z^{k+1}-z^k)\|^2.
\]
因此序列
\[
\Phi_k:=\|y^k-\bar y\|^2+\rho^2\|B(z^k-\bar z)\|^2
\]
单调不增且有下界，故收敛。于是
\[
\sum_{k=0}^\infty \|r^{k+1}\|^2<\infty,
\qquad
\sum_{k=0}^\infty \|B(z^{k+1}-z^k)\|^2<\infty,
\]
从而
\[
r^k\to 0.
\]
由定义~\ref{def:admm-iteration} 中对偶残差的定义，也得到
\[
s^k=\rho B^\ast(z^k-z^{k-1})\to 0.
\]

再取任意弱聚点 $(x^\infty,z^\infty,y^\infty)$。由 $r^k\to 0$，可得
\[
Ax^\infty+Bz^\infty=b.
\]
对 $x$-步与 $z$-步最优性条件取极限，并利用 $s^k\to 0$ 与图像闭性，得到
\[
0\in \partial f(x^\infty)+A^\ast y^\infty,
\qquad
0\in \partial g(z^\infty)+B^\ast y^\infty.
\]
这正是第 \ref{sec:7-4} 节定义~\ref{def:kkt-cone-system} 与推论~\ref{cor:kkt-via-fenchel-subdifferential} 在当前线性约束模型中的 KKT 条件。若 KKT 解集为单点，则所有聚点一致，因而整个序列弱收敛到该点。
\end{proof}

\begin{corollary}[ADMM 的 KKT 极限]\label{cor:admm-kkt-limit}
在定理~\ref{thm:admm-primal-dual-convergence} 的假设下，任意极限点 $(x^\infty,z^\infty,y^\infty)$ 都满足
\[
Ax^\infty+Bz^\infty=b,
\]
\[
0\in \partial f(x^\infty)+A^\ast y^\infty,
\qquad
0\in \partial g(z^\infty)+B^\ast y^\infty.
\]
因此，ADMM 的极限点正是相应原问题与对偶问题的 KKT 点。
\end{corollary}

\begin{proof}
这只是定理~\ref{thm:admm-primal-dual-convergence} 末尾结论的重述。
\end{proof}

\subsection{典型例子}

\begin{example}[两块变量的标准 ADMM]
当约束写成
\[
Lx-z=0
\]
时，ADMM 变成最常见的共识模型：$x$-步处理复杂损失，$z$-步处理简单正则或约束，乘子变量负责把两块重新绑在一起。把这个例子放在这里，是为了给第 \ref{sec:10-2} 节 DR 结构一个最直接的算法外观。
\end{example}

\begin{example}[函数空间上的一致性约束分裂]
在函数空间优化中，$x$ 可能代表状态变量，$z$ 代表控制或辅助变量，而 $Ax+Bz=b$ 编码 PDE、边界或一致性约束。此时 ADMM 的意义不在于矩阵分块，而在于把原本强耦合的 KKT 系统拆成两个局部可解子问题。
\end{example}

\begin{example}[推荐系统与分布式共识]
在推荐系统、联邦学习或分布式训练中，常把局部模型变量和全局共识变量分开建模，再用 ADMM 交替更新。把这个例子放在这里，是为了说明 Boyd 中“增广拉格朗日 + ADMM”的工程配方，本质上正是前七章抽象对象在算法层的一次计算实现。
\end{example}

\subsection*{有限维对照}

Boyd 对 ADMM 的经典叙述通常从增广拉格朗日和残差图表开始，这在有限维工程问题里非常有效。但如果把它只看成一个“经验上好用的迭代模板”，就会错过它与第 \ref{sec:10-2} 节 Douglas-Rachford 分裂、第 \ref{sec:7-4} 节 KKT 系统以及第 \ref{sec:9-4} 节单调包含之间的统一结构。本节的要点恰恰是：有限维矩阵公式只是抽象原--对偶分裂在 Euclidean 空间中的坍缩版本。

\subsection*{习题（待补）}

\begin{exercise}[增广拉格朗日占位]
补一题：对一个两块约束模型写出增广拉格朗日，并逐步推出定义~\ref{def:admm-iteration} 的三个更新式。
\end{exercise}

\begin{exercise}[DR--ADMM 接口题占位]
补一题：在一个简单共识模型中，说明 ADMM 如何由第 \ref{sec:10-2} 节的 Douglas--Rachford 分裂恢复出来。
\end{exercise}


\section{加速算法与 Hilbert 空间拓展}

\label{sec:10-4}

到这里为止，本章已经得到了两类稳定的 Hilbert 空间分裂模板：前向--后向分裂和 Douglas--Rachford/ADMM。下一步自然会问，能否像有限维凸优化那样再叠加一个“加速层”？答案并不是简单的肯定或否定。惰性项、Nesterov 型外推和 Lyapunov 能量确实能够在某些 Hilbert 空间模型中工作，但它们依赖的不是几句复杂度口号，而是比第 \ref{sec:10-1} 节更精细的能量耗散结构。也正因为如此，本节的重点不是罗列所有加速变体，而是给出一个可安全调用的模板，并说明其适用边界。

\subsection{惰性 proximal 模板}

\begin{definition}[惰性 proximal 迭代]\label{def:inertial-proximal-iteration}
设 $H$ 为 Hilbert 空间，考虑复合模型
\[
\min_{x\in H}\{h(x)+g(x)\},
\]
其中 $h$ 为凸且可微，$g$ 为 proper、凸且下半连续。给定步长 $\gamma>0$ 与惰性参数序列 $\{\alpha_k\}\subset [0,1)$，定义
\[
y^k:=x^k+\alpha_k(x^k-x^{k-1}),
\]
\[
x^{k+1}:=\operatorname{prox}_{\gamma g}\bigl(y^k-\gamma \nabla h(y^k)\bigr).
\]
该迭代称为\emph{惰性 proximal 迭代}。
\end{definition}

\begin{definition}[加速的 Lyapunov 能量]\label{def:lyapunov-energy-accelerated}
在定义~\ref{def:inertial-proximal-iteration} 的背景下，称形如
\[
\mathcal E_k:=\bigl(h(x^k)+g(x^k)\bigr)+\eta_k\|x^k-x^{k-1}\|^2
\]
的量为\emph{Lyapunov 能量}，其中 $\eta_k\ge 0$ 用于补偿惰性外推带来的动量项。
\end{definition}

\subsection{能量下降与模板性结论}

\begin{proposition}[加速迭代的能量下降]\label{prop:accelerated-energy-descent}
设 $h$ 的梯度是 $\beta$-cocoercive，$0\le \alpha_k\le \bar\alpha<1$，且步长 $\gamma$ 与惯性参数满足标准小步长条件。则可选择常数 $\eta>0$ 与 $c>0$，使定义~\ref{def:lyapunov-energy-accelerated} 中的能量满足
\[
\mathcal E_{k+1}+c\|x^{k+1}-x^k\|^2\le \mathcal E_k,
\qquad k\ge 0.
\]
\end{proposition}

\begin{proof}
由定义~\ref{def:inertial-proximal-iteration}，
\[
x^{k+1}=\operatorname{prox}_{\gamma g}(y^k-\gamma \nabla h(y^k)).
\]
这与第 \ref{sec:10-1} 节推论~\ref{cor:prox-gradient-specialization} 的 prox-gradient 步完全同型，只是把基点从 $x^k$ 换成了外推点 $y^k$。因此，对固定的 $y^k$，前向--后向步仍然给出一条下降不等式；而由第 \ref{sec:9-3} 节命题~\ref{prop:moreau-envelope-gradient} 所体现的平滑梯度结构，$\nabla h(y^k)$ 与 $\nabla h(x^k)$ 的差可以借助 cocoercive 性控制。

把 $y^k=x^k+\alpha_k(x^k-x^{k-1})$ 代入后，会多出若干交叉项。对这些项使用 Young 不等式，并把它们吸收到
\[
\eta\|x^k-x^{k-1}\|^2
\]
中，就可得到
\[
h(x^{k+1})+g(x^{k+1})
\le
h(x^k)+g(x^k)-c_1\|x^{k+1}-x^k\|^2+c_2\|x^k-x^{k-1}\|^2.
\]
只要步长和惰性参数满足标准小步长条件，就能选取 $\eta$ 使
\[
\mathcal E_{k+1}+c\|x^{k+1}-x^k\|^2\le \mathcal E_k
\]
成立。这里的关键不是常数的最优表达，而是存在一个真正耗散的 Lyapunov 能量，这正是惰性算法在 Hilbert 空间中可分析的根本原因。
\end{proof}

\begin{theorem}[Hilbert 空间中的加速模板]\label{thm:hilbert-acceleration-template}
设定义~\ref{def:inertial-proximal-iteration} 的问题存在解，且命题~\ref{prop:accelerated-energy-descent} 的能量下降成立。则：
\begin{enumerate}
  \item $\{\mathcal E_k\}$ 收敛，且
  \[
  \sum_{k=0}^{\infty}\|x^{k+1}-x^k\|^2<\infty;
  \]
  \item 每个弱聚点 $\bar x$ 都满足
  \[
  0\in \nabla h(\bar x)+\partial g(\bar x);
  \]
  \item 若解集退化为单点，或问题还满足额外的强凸/误差界条件，则整个序列弱收敛到该解。
\end{enumerate}
\end{theorem}

\begin{proof}
由命题~\ref{prop:accelerated-energy-descent}，$\{\mathcal E_k\}$ 单调不增且有下界，故收敛，并且
\[
\sum_{k=0}^{\infty}\|x^{k+1}-x^k\|^2<\infty.
\]
于是
\[
x^{k+1}-x^k\to 0,
\qquad
y^k-x^k=\alpha_k(x^k-x^{k-1})\to 0.
\]
再利用 prox 最优性条件，
\[
\frac{1}{\gamma}(y^k-x^{k+1})-\nabla h(y^k)\in \partial g(x^{k+1}).
\]
由于 $y^k-x^{k+1}\to 0$，而 $\nabla h$ 在有界集上连续，取弱聚点并使用次微分图像闭性，可得任意弱极限 $\bar x$ 满足
\[
-\nabla h(\bar x)\in \partial g(\bar x),
\]
也即
\[
0\in \nabla h(\bar x)+\partial g(\bar x).
\]
这正是第 \ref{sec:10-1} 节定理~\ref{thm:forward-backward-convergence} 所对应的零点条件。

若解集为单点，或额外具备强凸/误差界等消除多重聚点的结构，则 Opial 型论证可把“所有聚点都是解”升级为整列弱收敛。换言之，Hilbert 空间中的加速并不是单靠外推就自动成功，而是建立在命题~\ref{prop:accelerated-energy-descent} 的能量控制之上。
\end{proof}

\begin{remark}[无穷维中的加速边界]\label{remark:acceleration-limitations-infinite-dimensional}
有限维文献常把“加速”表述为一种几乎普适的复杂度提升，但在无穷维问题里，这样的说法往往过于乐观。首先，弱拓扑下缺乏自动紧性，意味着外推可能制造长时间振荡而不产生强收敛；其次，许多 Nesterov 型估计序列依赖 Euclidean 范数、有限维谱界或额外紧性，离开这些条件后就不能原封不动照搬。第 \ref{sec:11-1} 节会从连续时间角度再次看到这一点：真正稳定的加速模板必须伴随明确的 Lyapunov 能量，而不是只凭一句“更快”。
\end{remark}

\subsection{典型例子}

\begin{example}[惯性 proximal gradient]
在复合模型
\[
\min_x h(x)+g(x)
\]
中，若 $h$ 平滑、$g$ 可 prox，定义~\ref{def:inertial-proximal-iteration} 就给出了惯性 proximal gradient。它与第 \ref{sec:10-1} 节的前向--后向分裂只有一步之差：多出的外推项提高了经验效率，但也迫使我们引入定义~\ref{def:lyapunov-energy-accelerated} 的能量来控制稳定性。
\end{example}

\begin{example}[有限维 Nesterov 模板]
在 $\mathbb R^n$ 中，取 $g\equiv 0$ 时，惰性 proximal 迭代退化为带动量的梯度法。把这个例子放在这里，是为了提醒读者：Boyd 和经典一阶法文献中的加速直觉确实来自有限维，但本节的重点是说明它在 Hilbert 空间中仍可保留哪些结构、必须舍弃哪些口号。
\end{example}

\begin{example}[最优控制离散化后的加速求解]
在 PDE 约束或最优控制问题离散化后，常见的目标由光滑能量项和非光滑约束/正则项组成。惰性 proximal 更新往往能显著减少外层迭代数，但如果离散尺度变化导致条件数恶化，就必须重新检查命题~\ref{prop:accelerated-energy-descent} 的参数条件。把这个例子放在这里，是为了强调“加速”始终是受模型与空间结构制约的。
\end{example}

\subsection*{有限维对照}

Boyd 对加速算法的讨论通常以内点、梯度法或 Nesterov 迭代的复杂度估计为主，默认工作空间是 Euclidean 空间。本节的任务不是否认这些结论，而是把它们放回更准确的语境：在 Hilbert 空间里，能被安全继承的是 Lyapunov 能量、平均映射与零点结构；不能无条件继承的是所有基于有限维紧性和谱估计的复杂度口号。有限维结论仍然重要，但它们只是无穷维分析中的特例坍缩，而不是通用模板。

\subsection*{习题（待补）}

\begin{exercise}[Lyapunov 能量占位]
补一题：对一个具体的惯性 proximal gradient 迭代构造 Lyapunov 能量，并验证命题~\ref{prop:accelerated-energy-descent} 的下降格式。
\end{exercise}

\begin{exercise}[加速边界桥接题占位]
补一题：比较有限维 Nesterov 迭代和一个 Hilbert 空间惰性迭代，说明为什么后者必须额外检查 remark~\ref{remark:acceleration-limitations-infinite-dimensional} 中提到的稳定性边界。
\end{exercise}

